% Capítulo 3: Metodología
% Propósito: Detallar el enfoque metodológico, diseño de investigación, técnicas e instrumentos de recolección de información y la metodología de desarrollo de software seleccionada.

\section{Marco metodológico}
\label{sec:metodologia}

\subsection{Metodología de la investigación}

La presente investigación se orienta al diseño e implementación
de un sistema web de gestión de citas médicas con soporte para
la operación en entornos de conectividad intermitente. Para
organizar el proceso de investigación se propone la Metodología
de Investigación en Ciencia del Diseño
(\textit{Design Science Research Methodology}, DSRM), debido
a que permite estructurar el análisis de un problema, el diseño
y desarrollo de una solución tecnológica, y su posterior
demostración y evaluación. Esta metodología establece un proceso
orientado a la construcción de artefactos y a la valoración de
su utilidad para abordar un problema identificado.
\parencite{peffers2007dsrm}.

La metodología se utilizará como guía para relacionar la
problemática identificada con los requisitos del sistema, las
decisiones de diseño, la implementación del artefacto y la
comprobación de su funcionamiento. De esta manera, se busca
mantener la correspondencia entre los objetivos de la
investigación y los resultados obtenidos durante el desarrollo.

La aplicación de esta metodología se organizará en etapas
relacionadas con la identificación del problema, la definición
de los objetivos de la solución, el diseño y desarrollo del
sistema, su demostración, su evaluación y la comunicación de
los resultados.

\subsection{Tipo de investigación}

La investigación se orienta a la solución de un problema
práctico mediante el desarrollo de un artefacto de software.
Su propósito es diseñar e implementar un sistema web que
permita gestionar citas médicas y contemple mecanismos para
operar ante interrupciones de conectividad, de acuerdo con
los requisitos que se definan durante el análisis del sistema.

El resultado esperado es una implementación de software
acompañada de los modelos, diseños y documentación necesarios
para explicar su construcción y funcionamiento.

La delimitación definitiva del alcance de la investigación
y de los procedimientos de evaluación se establecerá en
correspondencia con los objetivos aprobados y los recursos
disponibles para el desarrollo del proyecto.

\subsection{Etapas metodológicas}

\subsection{Identificación del problema y motivación}

En esta etapa se delimitará el problema relacionado con la
gestión de citas médicas en condiciones de conectividad
intermitente. Se examinarán las necesidades que justifican
el desarrollo del sistema y las dificultades que pueden
presentarse cuando la conexión con los servicios remotos
no está disponible.

La identificación del problema servirá como fundamento para
establecer los requisitos y determinar qué operaciones
deberá contemplar el sistema durante una interrupción de
conectividad.

\subsection{Definición de los objetivos de la solución}

A partir del problema identificado y de los objetivos de
la investigación, se establecerán los requisitos funcionales
y no funcionales que deberá satisfacer el sistema.

Esta etapa permitirá delimitar las operaciones de gestión
de citas, los tipos de usuario contemplados y el comportamiento
esperado ante la pérdida y recuperación de la conexión.
Los requisitos definitivos se documentarán antes de utilizarlos
como criterios para el diseño y la evaluación del sistema.

\subsection{Diseño y desarrollo del sistema}

En esta etapa se elaborará el diseño de la solución tecnológica
a partir de los requisitos establecidos. Se definirán la
arquitectura del sistema, la organización de sus componentes,
los modelos de datos y los diagramas necesarios para
representar su estructura y funcionamiento.

También se diseñará el mecanismo que permitirá gestionar
las operaciones contempladas durante los periodos sin
conectividad y determinar cómo se procesarán los datos
pendientes cuando se recupere la conexión.

Posteriormente, se implementarán los componentes definidos
en el diseño. Las tecnologías y los mecanismos concretos
de almacenamiento local y sincronización se seleccionarán
de acuerdo con los requisitos y las condiciones técnicas
del proyecto.

\subsection{Demostración del sistema}

La demostración tendrá como finalidad comprobar que el
sistema implementado puede ejecutar las operaciones previstas
para los escenarios contemplados en los requisitos.

Se definirán casos de prueba relacionados con el flujo de
gestión de citas y con las condiciones de conectividad que
se decida evaluar. Los procedimientos definitivos dependerán
de la funcionalidad implementada y de las condiciones
disponibles para realizar las pruebas.

\subsection{Evaluación del sistema}

La evaluación se orientará a comprobar el cumplimiento de
los requisitos definidos y el comportamiento del sistema
en los escenarios de prueba establecidos.

En particular, se considerará la operación durante una
interrupción de conectividad y el procesamiento de los datos
pendientes tras la recuperación de la conexión, siempre
que estos comportamientos formen parte de los requisitos
implementados.

Los criterios, indicadores y procedimientos de evaluación
se especificarán a partir de los requisitos aprobados y de
las capacidades reales del sistema. La participación de
usuarios y la realización de pruebas en un entorno clínico
se determinarán de acuerdo con la disponibilidad de acceso
y las autorizaciones correspondientes.

\subsection{Comunicación de los resultados}

Finalmente, se documentarán el proceso de investigación,
las decisiones de diseño, la implementación realizada,
las pruebas ejecutadas y los resultados obtenidos.

La presentación de los resultados permitirá establecer la
correspondencia entre el problema identificado, los
objetivos de la investigación, la solución desarrollada
y las evidencias obtenidas durante su evaluación.



\clearpage

\section{Cronograma de actividades}
\label{sec:cronograma}

El cronograma presenta la distribución temporal de las
actividades previstas para el desarrollo del proyecto,
organizadas en dieciséis semanas de trabajo a partir
del 27 de julio de 2026.

\begin{figure}[htbp]
	\centering
	\scriptsize
	
	% Configuración general
	\setlength{\tabcolsep}{3pt}
	
	% =====================================================
	% PRIMER BLOQUE: SEMANAS 1 A 8
	% =====================================================
	
	\textbf{Primer período: semanas 1--8}
	
	\vspace{0.2cm}
	
	\begin{ganttchart}[
		hgrid,
		vgrid,
		x unit=0.68cm,
		y unit chart=0.48cm,
		title height=0.7,
		bar height=0.35,
		bar/.append style={fill=gray!45},
		bar label font=\tiny,
		bar label node/.append style={
			text width=5.2cm,
			align=left
		},
		title label font=\scriptsize\bfseries,
		title/.append style={fill=gray!20},
		canvas/.append style={draw=black},
		group label font=\tiny
		]{1}{8}
		
		\gantttitle{Semanas 1--8}{8}\\
		\gantttitlelist{1,...,8}{1}\\
		
		\ganttbar{1. Planificación y delimitación del problema}{1}{2}\\
		\ganttbar{2. Búsqueda de antecedentes y literatura científica}{2}{3}\\
		\ganttbar{3. Clasificación y análisis de la literatura}{3}{3}\\
		\ganttbar{4. Análisis y documentación de requisitos}{4}{5}\\
		\ganttbar{5. Historias de usuario y casos de uso}{5}{6}\\
		\ganttbar{6. Elaboración de diagramas UML}{6}{6}\\
		\ganttbar{7. Diseño de arquitectura y modelo de datos}{7}{7}\\
		\ganttbar{8. Diseño de operaciones sin conexión y preparación del entorno}{8}{8}
		
	\end{ganttchart}
	
	\vspace{0.6cm}
	
	% =====================================================
	% SEGUNDO BLOQUE: SEMANAS 9 A 16
	% =====================================================
	
	\textbf{Segundo período: semanas 9--16}
	
	\vspace{0.2cm}
	
	\begin{ganttchart}[
		hgrid,
		vgrid,
		x unit=0.68cm,
		y unit chart=0.48cm,
		title height=0.7,
		bar height=0.35,
		bar/.append style={fill=gray!45},
		bar label font=\tiny,
		bar label node/.append style={
			text width=5.2cm,
			align=left
		},
		title label font=\scriptsize\bfseries,
		title/.append style={fill=gray!20},
		canvas/.append style={draw=black},
		group label font=\tiny
		]{9}{16}
		
		\gantttitle{Semanas 9--16}{8}\\
		\gantttitlelist{9,...,16}{1}\\
		
		\ganttbar{9. Desarrollo de los módulos principales}{9}{10}\\
		\ganttbar{10. Implementación de operaciones sin conexión}{11}{11}\\
		\ganttbar{11. Implementación y ajustes de sincronización}{12}{12}\\
		\ganttbar{12. Pruebas funcionales y de conectividad}{13}{13}\\
		\ganttbar{13. Evaluación y corrección de errores}{14}{14}\\
		\ganttbar{14. Redacción de resultados y conclusiones}{15}{15}\\
		\ganttbar{15. Revisión y correcciones del documento}{16}{16}\\
		\ganttbar{16. Preparación de la entrega final}{16}{16}
		
	\end{ganttchart}
	
	\caption{Cronograma semanal previsto para el desarrollo del proyecto.}
	\label{fig:cronograma}
	
\end{figure}

\clearpage


